\SGANextSectionTarget{sga3:xxiii:subsection:3-2}\SGASubsection{3.2}{Groups of Type \(A_2\)}

\SGARefTarget{sga3:XXIII:3.2.1}
\SGAStatement{Proposition 3.2.1}{Let \(S\) be a scheme and \(G\) a
pinned \(S\)-group of type \(A_2\). Write
\[
  \Delta=\{\alpha,\beta\},
  \qquad
  R^+=\{\alpha,\beta,\alpha+\beta\}.
\]
\begin{enumerate}
  \item[\SGAInlineRefTarget{sga3:XXIII:3.2.1:item:i}\SGAInlineRefTarget{sga3:xxiii:proposition:3-2-1:item:i}(i)]
    \[
      t_{\alpha\beta}=(w_\alpha w_\beta)^3
        =e=(w_\beta w_\alpha)^3=t_{\beta\alpha}.
    \]
  \item[\SGAInlineRefTarget{sga3:xxiii:proposition:3-2-1:item:ii}\SGAInlineRefTarget{sga3:XXIII:3.2.1:item:ii}(ii)] Put
    \(\operatorname{Ad}(w_\beta)X_\alpha=X_{\alpha+\beta}\). Then
    \[
    \begin{aligned}
      \operatorname{Ad}(w_\alpha)X_\beta&=-X_{\alpha+\beta},\\
      \operatorname{Ad}(w_\alpha)X_{\alpha+\beta}&=X_\beta,\\
      \operatorname{Ad}(w_\beta)X_{\alpha+\beta}&=-X_\alpha.
    \end{aligned}
    \]
  \item[\SGAInlineRefTarget{sga3:XXIII:3.2.1:item:iii}(iii)] Put
    \[
      p_{\alpha+\beta}(x)
        =\exp(xX_{\alpha+\beta})
        =\operatorname{int}(w_\beta)p_\alpha(x).
    \]
    Then
    \[
      p_\beta(y)p_\alpha(x)
        =p_\alpha(x)p_\beta(y)p_{\alpha+\beta}(xy).
    \]
\end{enumerate}}

\SGARefTarget{sga3:XXIII:3.2.2}
\SGAParagraph{3.2.2}
The proof occupies \SGARefLink{sga3:XXIII:3.2.2}{3.2.2}--\SGARefLink{sga3:XXIII:3.2.7}{3.2.7}. First,
\[
  (\beta^*,\alpha)=(\alpha^*,\beta)=-1,
\]
so \(\alpha(t_\beta)=\beta(t_\alpha)=-1\). Put
\[
  \operatorname{Ad}(w_\beta)X_\alpha=X_{\alpha+\beta}.
\]
Then
\[
  \operatorname{Ad}(w_\beta)X_{\alpha+\beta}
    =\alpha(t_\beta)X_\alpha=-X_\alpha.
\]
Write
\[
  \operatorname{Ad}(w_\alpha)X_\beta=zX_{\alpha+\beta},
  \qquad z\in\mathbf G_m(S).
\]
It follows that
\[
  \operatorname{Ad}(w_\alpha)X_{\alpha+\beta}
    =\beta(t_\alpha)z^{-1}X_\beta=-z^{-1}X_\beta.
\]
By \SGARefLink{sga3:xxii:cor:5.5.2}{Exposé XXII, 5.5.2}, there is a unique
\(A\in\mathbf G_a(S)\) such that
\begin{equation}
  p_\beta(y)p_\alpha(x)
    =p_\alpha(x)p_\beta(y)p_{\alpha+\beta}(Axy).
  \tag{+}
  \label{sga3:XXIII:3.2:eqplus}
\end{equation}
Thus \SGARefLink{sga3:XXIII:3.2.1:item:ii}{(ii)} and \SGARefLink{sga3:XXIII:3.2.1:item:iii}{(iii)} reduce to proving \(A=-z=1\).

\SGARefTarget{sga3:XXIII:3.2.3}
\SGAParagraph{3.2.3}
Conjugating \eqref{sga3:XXIII:3.2:eqplus} by \(w_\beta\) gives
\begin{equation}
  p_{-\beta}(-y)p_{\alpha+\beta}(x)
    =p_{\alpha+\beta}(x)p_{-\beta}(-y)p_\alpha(-Axy).
  \tag{++}
  \label{sga3:XXIII:3.2:eqplusplus}
\end{equation}

\SGARefTarget{sga3:XXIII:3.2.4}
\SGAParagraph{3.2.4}
By definition,
\[
  w_\beta p_\alpha(x)w_\beta^{-1}=p_{\alpha+\beta}(x).
\]
Expanding \(w_\beta\), this is
\[
\begin{aligned}
  &p_\beta(1)p_{-\beta}(-1)p_\beta(1)p_\alpha(x)
    p_\beta(-1)p_{-\beta}(1)p_\beta(-1)\\
  &\hspace{55mm}=p_{\alpha+\beta}(x).
\end{aligned}
\]
Since \(p_\beta\) commutes with \(p_{\alpha+\beta}\), this may be
written
\[
  p_\beta(1)p_\alpha(x)p_\beta(-1)
    =p_{-\beta}(1)p_{\alpha+\beta}(x)p_{-\beta}(-1).
\]
Using \eqref{sga3:XXIII:3.2:eqplus} on the left and
\eqref{sga3:XXIII:3.2:eqplusplus} on the right gives
\[
\begin{aligned}
  p_\alpha(x)p_\beta(1)p_{\alpha+\beta}(Ax)p_\beta(-1)
   ={}&p_{\alpha+\beta}(x)p_{-\beta}(1)\\
     &{}\cdot p_\alpha(Ax)p_{-\beta}(-1).
\end{aligned}
\]
Because \(\alpha+2\beta\), \(\alpha-\beta\), and
\(2\alpha+\beta\) are not roots, this reduces to
\[
  p_\alpha(x)p_{\alpha+\beta}(Ax)
    =p_\alpha(Ax)p_{\alpha+\beta}(x).
\]
\SGARefLink{sga3:xxii:rem:4.1.3}{Exposé XXII, 4.1.3} now gives \(A=1\).

\SGARefTarget{sga3:XXIII:3.2.5}
\SGAParagraph{3.2.5}
Conjugating \eqref{sga3:XXIII:3.2:eqplus} by \(w_\alpha\), and using
\(A=1\), gives
\begin{equation}
  p_{\alpha+\beta}(zy)p_{-\alpha}(-x)
   =p_{-\alpha}(-x)p_{\alpha+\beta}(zy)
      p_\beta(-z^{-1}xy).
  \tag{+++}
  \label{sga3:XXIII:3.2:eqplusplusplus}
\end{equation}

\SGARefTarget{sga3:XXIII:3.2.6}
\SGAParagraph{3.2.6}
As before,
\[
  w_\alpha p_\beta(y)w_\alpha^{-1}=p_{\alpha+\beta}(zy).
\]
Since \(p_\alpha\) and \(p_{\alpha+\beta}\) commute,
\[
  p_\alpha(1)p_\beta(y)p_\alpha(-1)
    =p_{-\alpha}(1)p_{\alpha+\beta}(zy)p_{-\alpha}(-1).
\]
Using \eqref{sga3:XXIII:3.2:eqplus} and
\eqref{sga3:XXIII:3.2:eqplusplusplus}, this becomes
\[
  p_\beta(y)p_{\alpha+\beta}(-y)
    =p_{\alpha+\beta}(zy)p_\beta(-z^{-1}y).
\]
The two factors commute, and hence \(z=-1\).

\SGARefTarget{sga3:XXIII:3.2.7}
\SGAParagraph{3.2.7}
It remains to prove \SGARefLink{sga3:XXIII:3.2.1:item:i}{(i)}. One has
\[
  w_\alpha w_\beta w_\alpha
    =w_\alpha w_\beta w_\alpha^{-1}w_\alpha^2
    =w_{\alpha+\beta}^{-1}t_\alpha,
\]
and therefore
\[
\begin{aligned}
  w_\beta w_\alpha w_\beta w_\alpha w_\beta
    &=w_\beta w_{\alpha+\beta}^{-1}t_\alpha w_\beta\\
    &=w_\beta w_{\alpha+\beta}^{-1}w_\beta^{-1}
      \,s_\beta(t_\alpha)t_\beta\\
    &=w_\alpha t_\alpha t_\beta t_\beta
      =w_\alpha t_\alpha=w_\alpha^{-1}.
\end{aligned}
\]
Thus
\[
  (w_\alpha w_\beta)^3=(w_\beta w_\alpha)^3=e.
\]

\SGARefTarget{sga3:XXIII:3.2.8}
\SGAStatement{Remark 3.2.8}{Put
\(v_\alpha=f_\alpha(u_\alpha)\). Condition (v) of \SGARefLink{sga3:XXIII:2.3}{2.3} becomes
\[
\text{(A)}\quad
\begin{cases}
  \operatorname{int}(h_\alpha h_\beta h_\alpha)v_\beta=v_\beta,\\
  \operatorname{int}(h_\beta h_\alpha h_\beta)v_\alpha=v_\alpha,
\end{cases}
\]
and
\[
\text{(B)}\quad
\begin{cases}
  v_\beta v_\alpha
    =v_\alpha v_\beta\,(h_\beta v_\alpha h_\beta^{-1}),\\
  v_\beta(h_\beta v_\alpha h_\beta^{-1})
    =(h_\beta v_\alpha h_\beta^{-1})v_\beta,\\
  v_\alpha(h_\beta v_\alpha h_\beta^{-1})
    =(h_\beta v_\alpha h_\beta^{-1})v_\alpha.
\end{cases}
\]
Writing \(v_{\alpha+\beta}=\operatorname{int}(h_\beta)v_\alpha\),
the latter conditions are equivalently
\[
\text{(B)}\quad
\begin{cases}
  v_\beta v_\alpha=v_\alpha v_\beta v_{\alpha+\beta},\\
  v_\alpha v_{\alpha+\beta}=v_{\alpha+\beta}v_\alpha,\\
  v_\beta v_{\alpha+\beta}=v_{\alpha+\beta}v_\beta.
\end{cases}
\]}

\SGANextSectionTarget{sga3:xxiii:subsection:3-3}\SGASubsection{3.3}{Groups of Type \(B_2\)}

\SGARefTarget{sga3:XXIII:3.3.1}
\SGAInlineRefTarget{sga3:xxiii:proposition:3-3-1}\SGAStatement{Proposition 3.3.1}{Let \(S\) be a scheme and \(G\) a
pinned \(S\)-group of type \(B_2\). Write
\[
  \Delta=\{\alpha,\beta\},
  \qquad
  R^+=\{\alpha,\beta,\alpha+\beta,2\alpha+\beta\}.
\]
\begin{enumerate}
  \item[\SGAInlineRefTarget{sga3:xxiii:proposition:3-3-1:item:i}(i)]
    \[
      t_{\alpha\beta}=(w_\alpha w_\beta)^4
        =t_\alpha=(w_\beta w_\alpha)^4=t_{\beta\alpha}.
    \]
  \item[\SGAInlineRefTarget{sga3:XXIII:3.3.1:item:ii}\SGAInlineRefTarget{sga3:xxiii:proposition:3-3-1:item:ii}(ii)] If
    \[
      \operatorname{Ad}(w_\beta)X_\alpha=X_{\alpha+\beta},
      \qquad
      \operatorname{Ad}(w_\alpha)X_\beta=X_{2\alpha+\beta},
    \]
    then
    \[
    \begin{aligned}
      \operatorname{Ad}(w_\alpha)X_{\alpha+\beta}
        &=-X_{\alpha+\beta},&
      \operatorname{Ad}(w_\alpha)X_{2\alpha+\beta}&=X_\beta,\\
      \operatorname{Ad}(w_\beta)X_{\alpha+\beta}
        &=-X_\alpha,&
      \operatorname{Ad}(w_\beta)X_{2\alpha+\beta}&=X_{2\alpha+\beta}.
    \end{aligned}
    \]
  \item[\SGAInlineRefTarget{sga3:XXIII:3.3.1:item:iii}(iii)] Put
    \[
    \begin{aligned}
      p_{\alpha+\beta}(x)
        &=\exp(xX_{\alpha+\beta})
          =\operatorname{int}(w_\beta)p_\alpha(x),\\
      p_{2\alpha+\beta}(x)
        &=\exp(xX_{2\alpha+\beta})
          =\operatorname{int}(w_\alpha)p_\beta(x).
    \end{aligned}
    \]
    Then
    \[
    \begin{aligned}
      p_\beta(y)p_\alpha(x)
        &=p_\alpha(x)p_\beta(y)p_{\alpha+\beta}(xy)
          p_{2\alpha+\beta}(x^2y),\\
      p_{\alpha+\beta}(y)p_\alpha(x)
        &=p_\alpha(x)p_{\alpha+\beta}(y)
          p_{2\alpha+\beta}(2xy).
    \end{aligned}
    \]
\end{enumerate}}

\SGARefTarget{sga3:XXIII:3.3.2}
\SGAParagraph{3.3.2}
The proof occupies \SGARefLink{sga3:XXIII:3.3.2}{3.3.2}--\SGARefLink{sga3:XXIII:3.3.6}{3.3.6}. We have
\[
  (\beta^*,\alpha)=-1,\qquad(\alpha^*,\beta)=-2,
\]
so \(\alpha(t_\beta)=-1\) and \(\beta(t_\alpha)=1\). In particular,
\(\beta(t_\alpha)=\alpha(t_\alpha)=1\), so \(t_\alpha\) is a section
of \(\operatorname{Centr}(G)\). Put
\[
  \operatorname{Ad}(w_\beta)X_\alpha=X_{\alpha+\beta},
  \qquad
  \operatorname{Ad}(w_\alpha)X_\beta=X_{2\alpha+\beta}.
\]
Then
\[
\begin{aligned}
  \operatorname{Ad}(w_\beta)X_{\alpha+\beta}
    &=\alpha(t_\beta)X_\alpha=-X_\alpha,\\
  \operatorname{Ad}(w_\alpha)X_{2\alpha+\beta}
    &=\beta(t_\alpha)X_\beta=X_\beta,\\
  \operatorname{Ad}(w_\beta)X_{2\alpha+\beta}
    &=X_{2\alpha+\beta}.
\end{aligned}
\]
For some \(k\in\mathbf G_m(S)\),
\[
  \operatorname{Ad}(w_\alpha)X_{\alpha+\beta}
    =kX_{\alpha+\beta}.
\]
By \SGARefLink{sga3:xxii:cor:5.5.2}{Exposé XXII, 5.5.2}, there are
\(A,B,C\in\mathbf G_a(S)\) such that
\begin{align}
  p_\beta(y)p_\alpha(x)
    &=p_\alpha(x)p_\beta(y)p_{\alpha+\beta}(Axy)
      p_{2\alpha+\beta}(Bx^2y), \tag{1}\label{sga3:XXIII:3.3:eq1}\\
  p_{\alpha+\beta}(y)p_\alpha(x)
    &=p_\alpha(x)p_{\alpha+\beta}(y)
      p_{2\alpha+\beta}(Cxy). \tag{2}\label{sga3:XXIII:3.3:eq2}
\end{align}
It remains to prove
\[
  A=B=1,\qquad C=2,\qquad k=-1.
\]

\SGARefTarget{sga3:XXIII:3.3.3}
\SGAParagraph{3.3.3}
Conjugating \eqref{sga3:XXIII:3.3:eq1} by \(w_\alpha\) gives
\begin{equation}
\begin{aligned}
  p_{2\alpha+\beta}(y)p_{-\alpha}(-x)
    ={}&p_{-\alpha}(-x)p_{2\alpha+\beta}(y)\\
      &{}\cdot p_{\alpha+\beta}(Akxy)p_\beta(Bx^2y).
\end{aligned}
\tag{3}\label{sga3:XXIII:3.3:eq3}
\end{equation}
Transforming \eqref{sga3:XXIII:3.3:eq2} in the same way gives
\begin{equation}
  p_{\alpha+\beta}(ky)p_{-\alpha}(-x)
   =p_{-\alpha}(-x)p_{\alpha+\beta}(ky)p_\beta(Cxy).
\tag{4}\label{sga3:XXIII:3.3:eq4}
\end{equation}
Conjugating \eqref{sga3:XXIII:3.3:eq1} by \(w_\beta\) gives
\begin{equation}
\begin{aligned}
  p_{-\beta}(-y)p_{\alpha+\beta}(x)
    ={}&p_{\alpha+\beta}(x)p_{-\beta}(-y)\\
      &{}\cdot p_\alpha(-Axy)p_{2\alpha+\beta}(Bx^2y).
\end{aligned}
\tag{5}\label{sga3:XXIII:3.3:eq5}
\end{equation}

\SGARefTarget{sga3:XXIII:3.3.4}
\SGAParagraph{3.3.4}
Write
\[
  w_\beta p_\alpha(x)w_\beta^{-1}=p_{\alpha+\beta}(x).
\]
Because \(\alpha+2\beta\) is not a root,
\[
  p_\beta(1)p_\alpha(x)p_\beta(-1)
    =p_{-\beta}(1)p_{\alpha+\beta}(x)p_{-\beta}(-1).
\]
Using \eqref{sga3:XXIII:3.3:eq1} on the left and
\eqref{sga3:XXIII:3.3:eq5} on the right yields
\[
\begin{aligned}
  &p_\alpha(x)p_\beta(1)p_{\alpha+\beta}(Ax)
    p_{2\alpha+\beta}(Bx^2)p_\beta(-1)\\
  &\quad=p_{\alpha+\beta}(x)p_{-\beta}(1)p_\alpha(Ax)
    p_{2\alpha+\beta}(-Bx^2)p_{-\beta}(-1).
\end{aligned}
\]
After commuting the indicated factors and applying
\eqref{sga3:XXIII:3.3:eq2}, this becomes
\[
\begin{aligned}
  p_\alpha(x)p_{\alpha+\beta}(Ax)p_{2\alpha+\beta}(Bx^2)
  ={}&p_\alpha(Ax)p_{\alpha+\beta}(x)\\
     &{}\cdot p_{2\alpha+\beta}((AC-B)x^2).
\end{aligned}
\]
Hence
\[
  A=1,\qquad C=2B.
\]

\SGARefTarget{sga3:XXIII:3.3.5}
\SGAParagraph{3.3.5}
Now write
\[
  w_\alpha p_\beta(y)w_\alpha^{-1}=p_{2\alpha+\beta}(y).
\]
Since \(3\alpha+\beta\) is not a root,
\[
  p_\alpha(1)p_\beta(y)p_\alpha(-1)
    =p_{-\alpha}(1)p_{2\alpha+\beta}(y)p_{-\alpha}(-1).
\]
Using \eqref{sga3:XXIII:3.3:eq1} on the left and
\eqref{sga3:XXIII:3.3:eq3} on the right gives
\[
\begin{aligned}
  p_\beta(y)p_{\alpha+\beta}(-Ay)p_{2\alpha+\beta}(By)
  ={}&p_{2\alpha+\beta}(y)p_{\alpha+\beta}(Aky)p_\beta(By).
\end{aligned}
\]
The three groups commute, so
\[
  B=1,\qquad -A=Ak.
\]
Consequently
\[
  A=B=1,\qquad C=2,\qquad k=-1.
\]

\SGARefTarget{sga3:XXIII:3.3.6}
\SGAParagraph{3.3.6}
We have proved \SGARefLink{sga3:XXIII:3.3.1:item:ii}{(ii)} and \SGARefLink{sga3:XXIII:3.3.1:item:iii}{(iii)}. Since
\[
  s_\beta(t_\alpha)
    =t_\alpha t_\beta^{(\alpha^*,\beta)}
    =t_\alpha,
\]
one obtains successively
\[
\begin{aligned}
  w_\alpha w_\beta w_\alpha
    &=w_{2\alpha+\beta}t_\alpha,\\
  w_\beta w_\alpha w_\beta w_\alpha w_\beta
    &=w_{2\alpha+\beta}t_\alpha t_\beta,\\
  w_\alpha w_\beta w_\alpha w_\beta w_\alpha w_\beta w_\alpha
    &=w_\beta^{-1}t_\alpha.
\end{aligned}
\]
\NDE{The French prints
\((\alpha^*,\beta)=2\); the convention used immediately above gives
\(-2\). The conclusion is unchanged because \(t_\beta^2=e\).}
It follows that
\[
  (w_\beta w_\alpha)^4=t_\alpha,
  \qquad
  (w_\alpha w_\beta)^4=s_\beta(t_\alpha)=t_\alpha.
\]

\SGARefTarget{sga3:XXIII:3.3.7}
\SGAStatement{Remark 3.3.7}{Put
\[
  v_{\alpha+\beta}=\operatorname{int}(h_\beta)v_\alpha,
  \qquad
  v_{2\alpha+\beta}=\operatorname{int}(h_\alpha)v_\beta.
\]
Condition \SGARefLink{sga3:XXIII:2.3:item:v}{(v)} of \SGARefLink{sga3:XXIII:2.3}{2.3} becomes
\[
\text{(A)}\quad
\begin{cases}
  \operatorname{int}(h_\alpha h_\beta h_\alpha)v_\beta=v_\beta,\\
  \operatorname{int}(h_\beta h_\alpha h_\beta)v_\alpha=v_\alpha,
\end{cases}
\]
and
\[
\text{(B)}\quad
\begin{cases}
  v_\beta v_\alpha
    =v_\alpha v_\beta v_{\alpha+\beta}v_{2\alpha+\beta},\\
  v_{\alpha+\beta}v_\alpha
    =v_\alpha v_{\alpha+\beta}v_{2\alpha+\beta}^{\,2},\\
  v_{\alpha+\beta}v_\beta=v_\beta v_{\alpha+\beta},\\
  v_{2\alpha+\beta}v_\alpha=v_\alpha v_{2\alpha+\beta},\\
  v_{2\alpha+\beta}v_\beta=v_\beta v_{2\alpha+\beta},\\
  v_{2\alpha+\beta}v_{\alpha+\beta}
    =v_{\alpha+\beta}v_{2\alpha+\beta}.
\end{cases}
\]}
